% Worked Examples: Concept 2.5.3 - Laplace Transforms
\documentclass[11pt,a4paper]{article}
\usepackage{worked-examples-template}

\title{\textbf{Worked Examples}\\[0.5em]
\large Concept 2.5.3: Laplace Transforms\\[0.3em]
\normalsize \coursesubtitle{} -- \coursetitle}
\author{Assoc.\ Prof.\ William Robertson \and Dr Sean McGowan}
\date{\today}

\begin{document}

\maketitle
\tableofcontents
\newpage

\section{Concept 2.5.3: Laplace Transform}

\subsection{Example 1: Solving an ODE using Laplace Transforms}

\paragraph{Problem:} Use Laplace transforms to solve the differential equation:
\begin{align}
\ddot{y} + 5\dot{y} + 6y = 2u(t)
\end{align}

where $u(t)$ is a unit step function and initial conditions are $y(0) = 0$, $\dot{y}(0) = 0$.

\paragraph{Solution:}

Taking the Laplace transform of both sides (assuming zero initial conditions):
\begin{align}
s^2Y(s) + 5sY(s) + 6Y(s) = \frac{2}{s}
\end{align}

Note: $\mathcal{L}\{u(t)\} = \mathcal{L}\{1\} = \frac{1}{s}$ for a unit step.

Factor out $Y(s)$:
\begin{align}
Y(s)(s^2 + 5s + 6) = \frac{2}{s}
\end{align}

Solve for $Y(s)$:
\begin{align}
Y(s) = \frac{2}{s(s^2+5s+6)} = \frac{2}{s(s+2)(s+3)}
\end{align}

Use partial fraction decomposition:
\begin{align}
\frac{2}{s(s+2)(s+3)} = \frac{A}{s} + \frac{B}{s+2} + \frac{C}{s+3}
\end{align}

Multiply both sides by $s(s+2)(s+3)$:
\begin{align}
2 = A(s+2)(s+3) + Bs(s+3) + Cs(s+2)
\end{align}

Find coefficients:
\begin{itemize}
\item Set $s = 0$: $2 = A(2)(3) = 6A \Rightarrow A = \frac{1}{3}$
\item Set $s = -2$: $2 = B(-2)(-2+3) = -2B \Rightarrow B = -1$
\item Set $s = -3$: $2 = C(-3)(-3+2) = 3C \Rightarrow C = \frac{2}{3}$
\end{itemize}

Therefore:
\begin{align}
Y(s) = \frac{1/3}{s} + \frac{-1}{s+2} + \frac{2/3}{s+3}
\end{align}

Take the inverse Laplace transform using the table:
\begin{itemize}
\item $\mathcal{L}^{-1}\left\{\frac{1}{s}\right\} = 1$
\item $\mathcal{L}^{-1}\left\{\frac{1}{s+a}\right\} = e^{-at}$
\end{itemize}

Result:
\begin{align}
y(t) = \frac{1}{3} - e^{-2t} + \frac{2}{3}e^{-3t}, \quad t \geq 0
\end{align}

\textbf{Verification:} As $t \to \infty$, the exponential terms decay to zero:
\begin{align}
\lim_{t \to \infty} y(t) = \frac{1}{3}
\end{align}

This matches the expected steady-state value from the final value theorem:
\begin{align}
y_{ss} = \lim_{s \to 0} sY(s) = \lim_{s \to 0} s \cdot \frac{2}{s(s+2)(s+3)} = \frac{2}{2 \cdot 3} = \frac{1}{3}
\end{align}
\end{document}
